% ECONOMICS_PROBLEM: {"id": "J38-343", "title": "Minimum-wage thresholds and adjustment", "jel_code": "J38", "source_id": "OP-0351"}
% !TeX program = xelatex
\documentclass[11pt,a4paper]{article}
\usepackage[margin=23mm]{geometry}
\usepackage{fontspec}
\setmainfont{Latin Modern Roman}
\usepackage{amsmath,amssymb,mathtools}
\usepackage{xcolor,enumitem,booktabs,longtable,array,xurl,hyperref}
\definecolor{muted}{HTML}{53636C}
\hypersetup{colorlinks=true,urlcolor=blue,linkcolor=blue}
\setlength{\parindent}{0pt}
\setlength{\parskip}{5pt}
\newcommand{\R}{\mathbb R}
\newcommand{\E}{\mathbb E}
\newcommand{\N}{\mathbb N}
\newcommand{\ind}{\mathbf 1}
\newcommand{\Var}{\operatorname{Var}}
\newcommand{\Cov}{\operatorname{Cov}}
\newcommand{\supp}{\operatorname{supp}}
\DeclareMathOperator*{\argmin}{arg\,min}
\DeclareMathOperator*{\argmax}{arg\,max}
\newcommand{\entrymeta}[3]{{\sffamily\small\color{muted}\textbf{\texttt{#1}}\quad|\quad #2\par #3\par}}
\newcommand{\Problemref}[1]{\texttt{#1}}
\newcommand{\JELref}[2]{\texttt{#2} (JEL \texttt{#1})}
\begin{document}
\section*{Minimum-wage thresholds and adjustment}
\textbf{Problem ID:} \texttt{J38-343}\quad \textbf{JEL:} \texttt{J38}\par
% BEGIN ECONOMICS STATEMENT
\label{problem:OP-0351}
{\sffamily\small\color{muted}Original subject group: Labor markets and work\par}
\label{definitions:problem:30007633}
\textbf{Audit classification:} Published research agenda. Full text checked. The joint employment--price compatibility question is explicitly future work; a global minimum-wage threshold and automation extensions are editorial.
\subsection*{Definitions and precise research target}
\textbf{Complete editorial problem.} Fix a set of labor markets, a baseline wage floor \(m_0\), and a policy horizon \(T\). For each feasible minimum wage \(m\), let \(L_t(m)\), \(P_t(m)\), \(A_t(m)\), and \(Q_t(m)\) denote equilibrium employment, output prices, automation investment, and a predeclared job-quality index at date \(t\). These outcomes include employer entry, exit, and worker reallocation. Define \(\Delta_L(m,T)=E[L_T(m)-L_T(m_0)]\), with analogous changes for the other outcomes. Let \(\mathcal M\) be a specified class of models containing employer wage-setting power, heterogeneous product demand, and labor adjustment costs. The research task is to identify joint response curves over the observed policy support and determine whether one model in \(\mathcal M\) can explain their magnitudes simultaneously. If justified by the data, estimate \(m^*=\inf\{m>m_0:\Delta_L(m,T)<0\}\), defining it as infinity if the set is empty. This is a first-negative-response boundary, not proof of a unique crossing or permanently negative response above it. Do not assume a single threshold exists or extrapolate it across countries. Identification must state the policy variation, counterfactual trends, and treatment of spillovers. Distributional welfare conclusions additionally require explicit worker preferences, consumer incidence, and fiscal assumptions.

\textbf{Definitions and admissible scope.} A policy is a complete path \(m_t=m\) for \(t=0,\ldots,T\), with announcement date, enforcement, exceptions, and nominal-to-real conversion fixed. Let markets have baseline probability weights \(\mu_r\), summing to one. Define employment as paid worker-hours \(L_{rt}\), price as a fixed-basket index \(P_{rt}\), automation as real investment \(A_{rt}\), and quality as a bounded declared index \(Q_{rt}\). All are integrable potential outcomes under the whole policy path. Restrict \(m\) to a declared compact interval \(\mathcal D\) and set \(\Delta_L(m,T)=\sum_r\mu_rE[L_{rT}(m)-L_{rT}(m_0)]\). The infimum of the negative-response set is its first boundary; it is not necessarily a crossing point or a persistent adverse-effect threshold. Identification at each \(m\) requires support and consistency; counterfactuals outside policy support require structural restrictions. Supply the revenue, labor-supply, and adjustment laws defining \(\mathcal M\) before testing joint compatibility. Here \(T\) is a fixed nonnegative integer and \(\mathcal D\subseteq[m_0,\infty)\) is the declared feasible floor domain. In the boundary definition use \(\{m\in\mathcal D:m>m_0,\Delta_L(m,T)<0\}\). The data-law input contains market policy assignments, the four measured outcomes, and all conditioning covariates; the admissible structural family fixes exogeneity, trend, interference, demand, technology, entry, adjustment-cost, and equilibrium-selection restrictions. Joint compatibility means that the same member of this family generates the observed law and all four response curves; fitting separate models to separate outcomes does not suffice. Report identified sets at each feasible floor, and distinguish sampling intervals from structural identification intervals.
\subsection*{Variants and assumption changes}
\begin{enumerate}
\item Local version (editorial): on the observed interval \([m_0,m_0+h]\), identify the joint derivative \(\partial_mE[(L_T,P_T,A_T,Q_T)]\) under differentiability. This does not identify the global first-loss boundary.
\item Persistent-loss version (editorial): replace the first-negative boundary by \(\inf\{m\in\mathcal D:\Delta_L(u,T)<0\text{ for every }u\in\mathcal D,u\ge m\}\). Empty sets have value \(+\infty\). Compare short and long horizons, permitting nonmonotonicity and changing entry responses.
\end{enumerate}
\subsection*{References and source audit}
\textbf{Support and access.} Full text checked. The joint employment--price compatibility question is explicitly future work; a global minimum-wage threshold and automation extensions are editorial. \textbf{Locator:} Section 7, p. 714, immediately before Section 7.1.
\begin{enumerate}\raggedright
\item\hypertarget{definitions:source:30007633:1}{} Patrick Kline. 2025. \emph{Labor Market Monopsony: Fundamentals and Frontiers}. Minimum-wage discussion, printed p. 714, paragraph beginning with heterogeneous product demand.
\url{https://eml.berkeley.edu/~pkline/papers/HoLE_vol6_monopsony.pdf}.
DOI: \url{https://doi.org/10.1016/bs.heslab.2025.07.007}.
\end{enumerate}

\subsection*{Common conventions: Source mathematical conventions}

These conventions apply to each entry unless explicitly replaced.

\textbf{Models and identification.} A primitive model \(m\) specifies all probability laws, feasible choices, information, equilibrium and measurement rules used by its target. The admissible class \(\mathcal M\) is fixed independently of outcomes. If \(P_O(m)\) is its observable probability law and \(\tau(m)\) the target, its identified set at observable law \(P\) is
\[
 \mathcal I_\tau(P)=\{\tau(m):m\in\mathcal M,\ P_O(m)=P\}.
\]
Point identification means that this set is a singleton. An empty set signals incompatibility of data and assumptions. Coordinate bounds need not describe the joint set. Bounds are sharp only if every retained target is attainable or arbitrarily approachable by compatible models. Finite bounds require explicit restrictions on unobserved quantities; unrestricted confounding, hidden wealth, extrapolation or arbitrary equilibrium selection can yield uninformative sets.

\textbf{Causal interventions.} A potential outcome \(Y_i(p)\) is the outcome under a complete policy regime \(p\), including timing, financing, information and market spillovers. The target population and units are fixed before treatment. With interference, use the complete assignment vector or a justified exposure mapping; individual no-interference notation alone cannot identify an economy-wide intervention. A difference of mean potential outcomes does not require the cross-world joint distribution. Conditioning on post-treatment migration, survival, enrollment or employment changes the estimand and needs separate justification.

\textbf{Statistical statements.} An estimator, confidence set, forecast or test specifies a sampling law, model class and loss/coverage criterion. Uniform \((1-\alpha)\)-coverage means \(\inf_{m\in\mathcal M}\Pr_m\{\tau(m)\in C_n\}\geq1-\alpha\), or an explicitly asymptotic version. The whole parameter space trivially has coverage, so informativeness requires a stated width, risk or power objective. A forecasting improvement is relative to a named benchmark, information set and evaluation population.

\textbf{Equilibrium and optimization.} An equilibrium comprises feasible plans, optimality, beliefs where needed, budget constraints and all market/resource clearing. The primitives or a stated benchmark define each correspondence \(\mathcal E(p;m)\). Nonemptiness is assumed or proved, never inferred just from compact policy sets. A selected-equilibrium target uses a declared selection rule; a robust target quantifies over all permitted equilibria. Use suprema or infima unless attainment is established. An \(\varepsilon\)-optimal policy is feasible and within \(\varepsilon\) of the finite optimum value. Report an empty feasible set separately.

\textbf{Welfare and accounting.} Normative utility, interpersonal normalization, social weights, numeraire, discounting and the use of public receipts are primitives. Behavioral utility need not coincide with the welfare criterion. Household welfare includes ownership income and rents once; adding profits again to compensating variation generally double-counts them. A monetary equivalent is defined by an explicit transfer and price convention, with monotonicity and range conditions. Average effects, mechanism contributions, transfers and welfare losses are different targets. Infinite sums require convergence and dynamic feasible plans require terminal or no-Ponzi conditions.

\textbf{Source versions.} A working-paper date, revision date and journal publication year are distinguished. A DOI for a journal version is not automatically the DOI of an earlier manuscript. Locators refer to the stated version and printed pages where specified. A metadata-only check does not verify a claim about a full-text conclusion. Every entry's source subsection narrows the provenance claim accordingly.
% END ECONOMICS STATEMENT
\end{document}
